%% ----------------------------------------------------------------------------
%%  §1.11  Esercizi di riepilogo
%% ----------------------------------------------------------------------------
\section{Esercizi di riepilogo}
\label{sec:esercizi-cap1}

Gli esercizi di questo paragrafo non servono a «fare pratica» nel senso scolastico: ogni
soluzione è accompagnata da un commento sulle \emph{scelte logiche} che la guidano,
perché è lì che si impara a dimostrare. Il lettore provi sempre prima da solo.

\begin{esercizio}[Unicità di $\sqrt2$]
Dimostrare che esiste al più un numero reale positivo $c$ tale che $c^{2}=2$.
\begin{soluzione}
Siano $c_{1},c_{2}>0$ con $c_{1}^{2}=c_{2}^{2}=2$. Allora
$0=c_{1}^{2}-c_{2}^{2}=(c_{1}-c_{2})(c_{1}+c_{2})$; per la legge di annullamento del
prodotto (\cref{prop:annullamento}) uno dei fattori è nullo; ma $c_{1}+c_{2}>0$, quindi
$c_{1}=c_{2}$. \emph{Scelta logica:} per provare l'unicità si suppongono due oggetti e si
mostra che coincidono; la fattorizzazione $a^{2}-b^{2}=(a-b)(a+b)$ trasforma un
problema di uguaglianza in un problema di annullamento, dove abbiamo un teorema
pronto. Si noti che l'\emph{esistenza} di $c$ richiede invece l'assioma di completezza
(\cref{sec:nzq}): esistenza e unicità sono due questioni logicamente distinte.
\end{soluzione}
\end{esercizio}

\begin{esercizio}[Dove si usa la completezza?]
Si consideri l'insieme $A=\{x\in\Q : x<1\}$ e $B=\{x\in\Q : x\ge1\}$. Esiste un
razionale $c$ che separa $A$ e $B$? È in contraddizione con il
\cref{teo:Q-non-completo}?
\begin{soluzione}
Sì: $c=1$ separa $A$ e $B$ (anzi è l'unico separatore). Nessuna contraddizione: il
\cref{teo:Q-non-completo} afferma che l'assioma di completezza fallisce in $\Q$ per
\emph{almeno una} coppia di insiemi, non per tutte. \emph{Scelta logica:} la negazione di
«per ogni coppia $A,B$ esiste $c$» è «esiste una coppia $A,B$ per cui non esiste $c$»;
per confutare un'affermazione universale basta un controesempio, e per quello abbiamo
dovuto scegliere con cura gli insiemi della \eqref{eq:AB-radice2}.
\end{soluzione}
\end{esercizio}

\begin{esercizio}[Disuguaglianza tra media aritmetica e geometrica]
Dimostrare che per ogni $a,b\ge0$ si ha $\sqrt{ab}\le\dfrac{a+b}{2}$, con uguaglianza se
e solo se $a=b$.
\begin{soluzione}
Poiché entrambi i membri sono $\ge0$ e la funzione $t\mapsto t^{2}$ è strettamente
crescente su $[0,+\infty)$ (\cref{prop:potenza-crescente}), la disuguaglianza equivale
a $ab\le\frac{(a+b)^{2}}{4}$, cioè a $4ab\le a^{2}+2ab+b^{2}$, cioè a
$0\le a^{2}-2ab+b^{2}=(a-b)^{2}$, che è vera per \ref{ass:ordine-prodotto} (e la regola
dei segni). L'uguaglianza vale se e solo se $(a-b)^{2}=0$, cioè $a=b$. \emph{Scelta
logica:} la catena di «equivale a» è lecita solo perché a ogni passo abbiamo usato
un'equivalenza: la stretta crescenza del quadrato sui numeri non negativi, che vale nei
due versi, la moltiplicazione per $4>0$ (\cref{prop:ordine-prodotto}) e l'assioma
\ref{ass:ordine-somma}, che permette di sommare o sottrarre uno stesso termine ai due
membri; se avessimo usato una semplice implicazione, avremmo dovuto percorrere la
catena a ritroso.
\end{soluzione}
\end{esercizio}

\begin{marginfigure}
\centering
\begin{tikzpicture}[scale=0.9]
  \draw[grigio!70] (0,0) -- (4,0);
  \draw[grigio!70] (4,0) arc (0:180:2);
  \draw[blu, thick] (1.2,0) -- (1.2,{sqrt(1.2*2.8)});
  \draw[rame, thick] (2,0) -- (2,2);
  \draw[tratteggio] (0,0) -- (1.2,{sqrt(1.2*2.8)}) -- (4,0);
  \node[font=\scriptsize, below] at (0.6,0) {$a$};
  \node[font=\scriptsize, below] at (2.6,0) {$b$};
  \node[font=\scriptsize, blu, left] at (1.2,0.45) {$\sqrt{ab}$};
  \node[font=\scriptsize, rame, right] at (2,0.55) {$\tfrac{a+b}{2}$};
\end{tikzpicture}
\caption{Media aritmetica e geometrica: in una semicirconferenza di diametro $a+b$,
l'altezza $\sqrt{ab}$ (secondo teorema di Euclide) non supera mai il raggio.}
\end{marginfigure}

\begin{esercizio}[Funzioni e monotonia]
Sia $f\colon\R\to\R$ strettamente crescente. Dimostrare che per ogni $y\in\R$
l'equazione $f(x)=y$ ha \emph{al più} una soluzione. Mostrare con un esempio che può
non averne nessuna.
\begin{soluzione}
Se $f(x_{1})=f(x_{2})=y$ con $x_{1}\neq x_{2}$, diciamo $x_{1}<x_{2}$, la stretta
crescenza darebbe $f(x_{1})<f(x_{2})$: assurdo. Per l'esempio: $f(x)=e^{x}$ è
strettamente crescente, ma $e^{x}=-1$ non ha soluzioni. \emph{Scelta logica:} «al più
una» è un'affermazione di unicità e si tratta come sempre; l'esempio mostra che
l'unicità non implica l'esistenza, cioè che iniettività e suriettività sono proprietà
indipendenti (\cref{fig:iniettiva-suriettiva}).
\end{soluzione}
\end{esercizio}

\begin{esercizio}[Valore assoluto]
Dimostrare che per ogni $x,y\in\R$ vale $\abs{x\,y}=\abs{x}\,\abs{y}$ \emph{senza}
distinguere i casi, usando l'identità $\abs{t}=\sqrt{t^{2}}$.
\begin{soluzione}
Si ha $\abs{t}^{2}=t^{2}$ per ogni $t$ (in entrambi i casi della definizione), e
$\abs{t}\ge0$; poiché la radice quadrata è l'inversa del quadrato su $[0,+\infty)$,
$\abs{t}=\sqrt{t^{2}}$. Allora
$\abs{xy}=\sqrt{(xy)^{2}}=\sqrt{x^{2}y^{2}}=\sqrt{x^{2}}\sqrt{y^{2}}=\abs{x}\abs{y}$,
dove abbiamo usato $\sqrt{uv}=\sqrt u\sqrt v$ per $u,v\ge0$ (che segue dall'unicità
della radice: $(\sqrt u\sqrt v)^{2}=uv$). \emph{Scelta logica:} rappresentare un
oggetto definito per casi con una formula unica elimina la distinzione dei casi; il
prezzo è dover usare proprietà della radice, cioè in ultima analisi la completezza.
Entrambe le dimostrazioni sono corrette: quella per casi usa meno strumenti, questa è
più corta.
\end{soluzione}
\end{esercizio}

\begin{esercizio}[Induzione: dove sta l'errore?]
Si consideri la seguente «dimostrazione» che tutti i cavalli hanno lo stesso colore.
\emph{Per $n=1$ è ovvio. Supponiamo che in ogni insieme di $n$ cavalli tutti abbiano lo
stesso colore, e consideriamo $n+1$ cavalli $c_{1},\dots,c_{n+1}$. I primi $n$ hanno lo
stesso colore per ipotesi induttiva; anche gli ultimi $n$, $c_{2},\dots,c_{n+1}$, hanno lo
stesso colore. Poiché $c_{2}$ sta in entrambi i gruppi, tutti gli $n+1$ cavalli hanno lo
stesso colore.} Dove sta l'errore?
\begin{soluzione}
Il passo induttivo fallisce per $n=1$: i due gruppi $\{c_{1}\}$ e $\{c_{2}\}$ non hanno
elementi in comune, e il «ponte» $c_{2}$ non esiste. Il passo $\mathcal{P}(n)\implica
\mathcal{P}(n+1)$ è corretto per ogni $n\ge2$, ma la base $\mathcal{P}(1)$ non si
collega a $\mathcal{P}(2)$: la prima tessera del domino cade, ma non urta la seconda.
\emph{Scelta logica:} nel verificare il passo induttivo bisogna assicurarsi che il
ragionamento valga \emph{per ogni} $n$ a partire dalla base, inclusi i primi valori,
che sono spesso quelli critici.
\end{soluzione}
\end{esercizio}

\needspace{14\baselineskip}
\begin{esercizio}[Con Python]
Scrivere una funzione Python \texttt{radice(y, n)} che, dato $y\ge0$, restituisca
un'approssimazione di $\sqrt{y}$ con errore minore di $10^{-n}$ usando il metodo di
bisezione del \cref{sec:nzq}: si parte dall'intervallo $[0,\max\{1,y\}]$ e lo si dimezza
finché la sua lunghezza è minore di $10^{-n}$. Quante iterazioni servono per $n=12$?
Perché il metodo funziona (quale teorema di questo capitolo garantisce che il punto
medio $m$ soddisfa $m^{2}<y$ oppure $m^{2}\ge y$, e che la scelta è coerente con
l'ordine)?
\begin{soluzione}
\begin{python}
def radice(y, n):
    a, b = 0.0, max(1.0, y)          # a^2 <= y <= b^2
    iterazioni = 0
    while b - a >= 10**(-n):
        m = (a + b) / 2
        if m * m < y:
            a = m                    # m è "troppo piccolo"
        else:
            b = m                    # m è "troppo grande"
        iterazioni += 1
    return (a + b) / 2, iterazioni

print(radice(2, 12))                 # (1.414213562372879, 41)
\end{python}
La dicotomia \ref{ass:dicotomia} garantisce che vale $m^{2}<y$ oppure $m^{2}\ge y$;
la \cref{prop:potenza-crescente} (stretta crescenza di $t\mapsto t^{2}$ sui non
negativi) garantisce che se $m^{2}<y$ allora la radice cercata è maggiore di $m$, e
viceversa. Servono $41$ iterazioni: l'intervallo iniziale $[0,2]$ ha lunghezza $2$, dopo
$k$ dimezzamenti ha lunghezza $2^{1-k}$, e $2^{-40}\approx9{,}1\cdot10^{-13}$ è la prima
potenza sotto $10^{-12}$. Ogni passo guadagna una cifra binaria, cioè circa $0{,}3$
cifre decimali. \emph{Scelta logica:} la bisezione usa soltanto l'ordinamento, e
l'invariante $a^{2}\le y\le b^{2}$, conservato a ogni passo, garantisce che la radice
cercata stia sempre in $[a,b]$: l'errore è quindi controllato a priori dalla lunghezza
dell'intervallo, che si dimezza, e il numero di iterazioni si calcola in anticipo.
\end{soluzione}
\end{esercizio}
